% GENERATED FILE. Edit canonical lessons, then run npm run generate:books.
% canonical-source-hash: fnv1a64:6b908da8edf6bfb5
% canonical-lessons: ML-C386-atitheyan, ML-C386-atithyam, ML-C386-svikaranam, ML-C386-ksanikkuka, ML-C386-svikarikkuka

\chapter{Hosts and Guests}
\label{ch:ml-hosts-and-guests}
\hlchaptermodality{386}

\begin{chapteropening}
\textbf{By the end of this chapter:} \emph{I can say a host, hospitality, a reception, to invite and to receive.}

Complete the last lesson of chapter 386: I can say a host, hospitality, a reception, to invite and to receive.
\end{chapteropening}

\section[\texorpdfstring{\ml{ആതിഥേയൻ} (ātithēyan)}{ആതിഥേയൻ (ātithēyan)}]{\ml{ആതിഥേയൻ} (ātithēyan) — a host}

\label{lesson:ML-C386-atitheyan}



\begin{quote}
Before the new one: say the Malayalam for an invitation card, then the Malayalam for to marry.
\end{quote}

\subsection*{You'll want to know: \ml{ആതിഥേയൻ}}
\textbf{\ml{ആതിഥേയൻ}} — \emph{ātithēyan} — \textquotedblleft{}a host\textquotedblright{}.

\begin{quote}
\ml{ആതിഥേയൻ} \ml{എല്ലാ} \ml{അതിഥികളെയും} \ml{വാതിൽക്കൽ} \ml{സ്വാഗതം} \ml{ചെയ്തു}.
\end{quote}

\begin{quote}
\emph{ātithēyan ellā atithikaḷeyuṁ vātilkkal svāgataṁ ceytu.}
\end{quote}

\begin{quote}
The host welcomed all the guests at the door.
\end{quote}

From \ml{അതിഥി}, a guest: the one who receives guests.

\begin{cousinweb}
\textbf{\ml{ആതിഥേയൻ}} (\emph{ātithēyan}) is borrowed from Sanskrit \emph{ātitheya}, \textquotedblleft{}proper for a guest, hospitable\textquotedblright{} (Monier-Williams), from \emph{atithi}, \textquotedblleft{}a guest\textquotedblright{}, the source of \ml{അതിഥി} (\emph{atithi}).
\end{cousinweb}

\begin{grammarlens}[title={What you've built}]
One more word for telling what happened.
\end{grammarlens}

\subsection*{Your turn}
\begin{itemize}
\raggedright
  \item \emph{Say it:} \emph{ātithēyan}
  \item \emph{Say it:} \emph{ātithēyan}, once more
  \item \emph{Say it:} \emph{ātithēyan ellā atithikaḷeyuṁ vātilkkal svāgataṁ ceytu.} — The host welcomed all the guests at the door.
  \item \emph{Recall:} say the Malayalam for an invitation card, then the Malayalam for to marry, then say \emph{ātithēyan} again
\end{itemize}

\begin{tcolorbox}[breakable,colback=teal!4,colframe=teal!35!black,title={Before you move on}]
What does \ml{ആതിഥേയൻ} mean? (\textquotedblleft{}A host\textquotedblright{}.) Say it once more.
\end{tcolorbox}
\section[\texorpdfstring{\ml{ആതിഥ്യം} (ātithyaṁ)}{ആതിഥ്യം (ātithyaṁ)}]{\ml{ആതിഥ്യം} (ātithyaṁ) — hospitality}

\label{lesson:ML-C386-atithyam}



\begin{quote}
Before the new one: say the Malayalam for to marry, then the Malayalam for a host.
\end{quote}

\subsection*{You'll want to know: \ml{ആതിഥ്യം}}
\textbf{\ml{ആതിഥ്യം}} — \emph{ātithyaṁ} — \textquotedblleft{}hospitality\textquotedblright{}.

\begin{quote}
\ml{അവരുടെ} \ml{ആതിഥ്യം} \ml{ഞങ്ങൾ} \ml{ഒരിക്കലും} \ml{മറക്കില്ല}.
\end{quote}

\begin{quote}
\emph{avaruṭe ātithyaṁ ñaṅṅaḷ orikkaluṁ maṟakkilla.}
\end{quote}

\begin{quote}
We will never forget their hospitality.
\end{quote}

Also from \ml{അതിഥി}: what a good host gives.

\begin{cousinweb}
\textbf{\ml{ആതിഥ്യം}} (\emph{ātithyaṁ}) is borrowed from Sanskrit \emph{ātithya}, \textquotedblleft{}proper for a guest\textquotedblright{} (Monier-Williams), the noun of \emph{atithi}, \textquotedblleft{}a guest\textquotedblright{}.
\end{cousinweb}

\begin{grammarlens}[title={What you've built}]
One more word for telling what happened.
\end{grammarlens}

\subsection*{Your turn}
\begin{itemize}
\raggedright
  \item \emph{Say it:} \emph{ātithyaṁ}
  \item \emph{Say it:} \emph{ātithyaṁ}, once more
  \item \emph{Say it:} \emph{avaruṭe ātithyaṁ ñaṅṅaḷ orikkaluṁ maṟakkilla.} — We will never forget their hospitality.
  \item \emph{Recall:} say the Malayalam for to marry, then the Malayalam for a host, then say \emph{ātithyaṁ} again
\end{itemize}

\begin{tcolorbox}[breakable,colback=teal!4,colframe=teal!35!black,title={Before you move on}]
What does \ml{ആതിഥ്യം} mean? (\textquotedblleft{}Hospitality\textquotedblright{}.) Say it once more.
\end{tcolorbox}
\section[\texorpdfstring{\ml{സ്വീകരണം} (svīkaraṇaṁ)}{സ്വീകരണം (svīkaraṇaṁ)}]{\ml{സ്വീകരണം} (svīkaraṇaṁ) — a reception; a welcome}

\label{lesson:ML-C386-svikaranam}



\begin{quote}
Before the new one: say the Malayalam for a host, then the Malayalam for hospitality.
\end{quote}

\subsection*{You'll want to know: \ml{സ്വീകരണം}}
\textbf{\ml{സ്വീകരണം}} — \emph{svīkaraṇaṁ} — \textquotedblleft{}a reception; a welcome\textquotedblright{}.

\begin{quote}
\ml{വിവാഹത്തിന്} \ml{ശേഷം} \ml{ഹാളിൽ} \ml{ഒരു} \ml{സ്വീകരണം} \ml{ഉണ്ടായിരുന്നു}.
\end{quote}

\begin{quote}
\emph{vivāhattinŭ śēṣaṁ hāḷil oru svīkaraṇaṁ uṇṭāyirunnu.}
\end{quote}

\begin{quote}
After the wedding there was a reception in the hall.
\end{quote}

Also a welcome given to someone: \ml{അവന്} \ml{വലിയ} \ml{സ്വീകരണം} \ml{കിട്ടി}.

\begin{cousinweb}
\textbf{\ml{സ്വീകരണം}} (\emph{svīkaraṇaṁ}) is borrowed from Sanskrit \emph{svīkaraṇa}, \textquotedblleft{}making one's own, accepting\textquotedblright{} (Monier-Williams): a reception takes guests in as one's own.
\end{cousinweb}

\begin{grammarlens}[title={What you've built}]
One more word for telling what happened.
\end{grammarlens}

\subsection*{Your turn}
\begin{itemize}
\raggedright
  \item \emph{Say it:} \emph{svīkaraṇaṁ}
  \item \emph{Say it:} \emph{svīkaraṇaṁ}, once more
  \item \emph{Say it:} \emph{vivāhattinŭ śēṣaṁ hāḷil oru svīkaraṇaṁ uṇṭāyirunnu.} — After the wedding there was a reception in the hall.
  \item \emph{Recall:} say the Malayalam for a host, then the Malayalam for hospitality, then say \emph{svīkaraṇaṁ} again
\end{itemize}

\begin{tcolorbox}[breakable,colback=teal!4,colframe=teal!35!black,title={Before you move on}]
What does \ml{സ്വീകരണം} mean? (\textquotedblleft{}A reception; a welcome\textquotedblright{}.) Say it once more.
\end{tcolorbox}
\section[\texorpdfstring{\ml{ക്ഷണിക്കുക} (kṣaṇikkuka)}{ക്ഷണിക്കുക (kṣaṇikkuka)}]{\ml{ക്ഷണിക്കുക} (kṣaṇikkuka) — to invite}

\label{lesson:ML-C386-ksanikkuka}



\begin{quote}
Before the new one: say the Malayalam for hospitality, then the Malayalam for a reception.
\end{quote}

\subsection*{You'll want to know: \ml{ക്ഷണിക്കുക}}
\textbf{\ml{ക്ഷണിക്കുക}} — \emph{kṣaṇikkuka} — \textquotedblleft{}to invite\textquotedblright{}.

\begin{quote}
\ml{അവർ} \ml{ഞങ്ങളെ} \ml{വിവാഹത്തിന്} \ml{ക്ഷണിച്ചു}.
\end{quote}

\begin{quote}
\emph{avar ñaṅṅaḷe vivāhattinŭ kṣaṇiccu.}
\end{quote}

\begin{quote}
They invited us to the wedding.
\end{quote}

Past: \ml{ക്ഷണിച്ചു}, \emph{kṣaṇiccu}.

The guest takes -\ml{എ}, the event -\ml{ക്ക്} or -\ml{ന്}: \ml{എന്നെ} \ml{പാർട്ടിക്ക്} \ml{ക്ഷണിച്ചു}.

\begin{cousinweb}
\textbf{\ml{ക്ഷണിക്കുക}} (\emph{kṣaṇikkuka}) is built on \ml{ക്ഷണം} (\emph{kṣaṇaṁ}), \textquotedblleft{}an invitation\textquotedblright{}, from Sanskrit \emph{kṣaṇa} (CDIAL 3642), with the Malayalam verb ending -\ml{ിക്കുക}.
\end{cousinweb}

\begin{grammarlens}[title={What you've built}]
One more word for telling what happened.
\end{grammarlens}

\subsection*{Your turn}
\begin{itemize}
\raggedright
  \item \emph{Say it:} \emph{kṣaṇikkuka}
  \item \emph{Say it:} \emph{kṣaṇikkuka}, once more
  \item \emph{Say it:} \emph{avar ñaṅṅaḷe vivāhattinŭ kṣaṇiccu.} — They invited us to the wedding.
  \item \emph{Recall:} say the Malayalam for hospitality, then the Malayalam for a reception, then say \emph{kṣaṇikkuka} again
\end{itemize}

\begin{tcolorbox}[breakable,colback=teal!4,colframe=teal!35!black,title={Before you move on}]
What does \ml{ക്ഷണിക്കുക} mean? (\textquotedblleft{}To invite\textquotedblright{}.) Say it once more.
\end{tcolorbox}
\section[\texorpdfstring{\ml{സ്വീകരിക്കുക} (svīkarikkuka)}{സ്വീകരിക്കുക (svīkarikkuka)}]{\ml{സ്വീകരിക്കുക} (svīkarikkuka) — to receive; to accept}

\label{lesson:ML-C386-svikarikkuka}



\begin{quote}
Before the new one: say the Malayalam for a reception, then the Malayalam for to invite.
\end{quote}

\subsection*{You'll want to know: \ml{സ്വീകരിക്കുക}}
\textbf{\ml{സ്വീകരിക്കുക}} — \emph{svīkarikkuka} — \textquotedblleft{}to receive; to accept\textquotedblright{}.

\begin{quote}
\ml{അമ്മ} \ml{അതിഥികളെ} \ml{ചിരിയോടെ} \ml{സ്വീകരിച്ചു}.
\end{quote}

\begin{quote}
\emph{amma atithikaḷe ciriyōṭe svīkariccu.}
\end{quote}

\begin{quote}
Mother received the guests with a smile.
\end{quote}

Past: \ml{സ്വീകരിച്ചു}, \emph{svīkariccu}.

Also to accept: \ml{സമ്മാനം} \ml{സ്വീകരിച്ചു}, accepted the prize.

\begin{cousinweb}
\textbf{\ml{സ്വീകരിക്കുക}} (\emph{svīkarikkuka}) is built from Sanskrit \emph{svīkāra}, \textquotedblleft{}making one's own; reception\textquotedblright{} (Monier-Williams), with the Malayalam verb ending -\ml{ിക്കുക}.
\end{cousinweb}

\begin{grammarlens}[title={What you've built}]
One more word for telling what happened.
\end{grammarlens}

\subsection*{Your turn}
\begin{itemize}
\raggedright
  \item \emph{Say it:} \emph{svīkarikkuka}
  \item \emph{Say it:} \emph{svīkarikkuka}, once more
  \item \emph{Say it:} \emph{amma atithikaḷe ciriyōṭe svīkariccu.} — Mother received the guests with a smile.
  \item \emph{Recall:} say the Malayalam for a reception, then the Malayalam for to invite, then say \emph{svīkarikkuka} again
\end{itemize}

\begin{tcolorbox}[breakable,colback=teal!4,colframe=teal!35!black,title={Before you move on}]
What does \ml{സ്വീകരിക്കുക} mean? (\textquotedblleft{}To receive; to accept\textquotedblright{}.) Say it once more.
\end{tcolorbox}
